% !PRINTED 477
\weberpage{477}
If one applies to this the substitutions $\omega$ and $\tau$, one sees that $v$ remains unchanged and hence admits all permutations of the group $P_{168}$.

Now $P_{168}$ is a divisor of the symmetric permutation group of seven elements of index 30, and a divisor of the alternating group of index 15. Consequently $v$ is a root of an equation of the 30th degree, whose coefficients are symmetric functions of the $\lambda$, and a root of an equation of the 15th degree, whose coefficients also contain the difference product of the $\lambda$. If the $\lambda$ are the roots of an equation of the 7th degree without affect, then $v$ is the root of a resolvent of the 30th degree, which by adjunction of the square root of the discriminant splits into two factors of the 15th degree.

If, besides the symmetric functions of the $\lambda$, the quantity $v$ belongs to the domain of rationality --- whether it is rational from the outset, or the domain of rationality is extended by the adjunction of $v$ --- then the $\lambda$ are the roots of a special equation of the 7th degree whose \gesperrt{Galois} group is of the 168th degree. What we now still have to prove is that this special kind of equations of the 7th degree can be reduced to the form problem of the group $G_{168}$.

To achieve this, we must seek to form three rational functions $X_1$, $X_2$, $X_3$ of the roots $\lambda$ which, when the permutations of the group $P_{168}$ are carried out, undergo the corresponding linear substitutions of the group $G_{168}$, i.\,e. which, if $\pi$ is a permutation from $P_{168}$ and $A$ is the corresponding substitution from $G_{168}$, pass into $A(X_1, X_2, X_3)$ upon carrying out the permutation $\pi$.

If we substitute these functions $X_1, X_2, X_3$ for the variables in the invariants of the group $G_{168}$, these invariants pass into functions of the $\lambda_i$ which remain unchanged under the permutations of the group $P_{168}$, and which consequently belong to the domain of rationality. The computation of the values of the functions $X_1, X_2, X_3$ from these values of the invariants is then the form problem for the $G_{168}$. Any function of the $X_1, X_2, X_3$ that is changed by the substitutions of $G_{168}$ is a root of a resolvent of the given equation of the 7th degree, and indeed, since the groups $G_{168}$ and $P_{168}$ are simple, a total resolvent.
